\section{Method}\label{sec:method}
% Objective: Quotes to forwards to SVI under arbitrage constraints, Dupire local volatility, correlated multi-asset simulation, Longstaff-Schwartz for the issuer's call, and the funding-spread inversion.

\subsection{The volatility surface}

Each expiry is fitted separately in total implied variance against
log-moneyness, $\totvar(\logm) = \sigma_{\mathrm{BS}}^2(\logm)\,T$ with
$\logm = \ln(K/F_T)$, using the raw SVI parameterisation of
\citet{GatheralJacquier2014},
%
\begin{equation}
  \totvar(\logm) = a + b\left[\rho\,(\logm - m)
      + \sqrt{(\logm - m)^2 + s^2}\,\right].
  \label{eq:svi}
\end{equation}
%
Fitting variance rather than volatility is a choice the next step forces. Local
volatility needs $\partial\totvar/\partial T$ and
$\partial^2\totvar/\partial\logm^2$, and \eqref{eq:svi} is twice differentiable
in $\logm$ in closed form, so those derivatives come off a fitted slice instead
of off differenced quotes.

Two arbitrage conditions are imposed, both from \citet{GatheralJacquier2014}.
The risk-neutral density implied by a slice is non-negative exactly when
%
\begin{equation}
  g(\logm) = \left(1 - \frac{\logm\,\totvar'}{2\totvar}\right)^{\!2}
    - \frac{(\totvar')^2}{4}\left(\frac{1}{\totvar} + \frac14\right)
    + \frac{\totvar''}{2} \;\geq\; 0
  \label{eq:butterfly}
\end{equation}
%
everywhere, and total variance must not decrease in $T$ at fixed $\logm$, or a
calendar spread prices negative. Neither check is cosmetic. A slice violating
\eqref{eq:butterfly} has negative density somewhere, and the denominator of the
local-variance expression below passes through zero at that point, which is the
usual reason a local-volatility surface comes back full of undefined values.
When a check fails the slice is refitted subject to the constraint rather than
repaired afterwards, the arbitrage-free parameter domain being characterised by
\citet{MartiniMingone2022}.

Where the conditions are enforced turns out to matter as much as whether they
are. Checking on one fixed wide grid asks a two-week slice about strikes far
below its deepest quote, where it is purely extrapolating, and enforcing the
calendar condition out there makes each short-dated slice's extrapolated
wing a floor for the next expiry, so the constraint propagates forward from the
slice carrying the least information rather than the most. The grid is scoped
to the strikes each pair of adjacent expiries actually quotes, padded slightly
and capped at five standard deviations of total volatility. Quotes enter the fit
only if the expiry is at least a fortnight out and the relative bid-ask is at
most 0.35, with log-moneyness clipped to $[-0.9, 0.45]$, and consecutive
expiries closer than 25 days apart are thinned so that the calendar condition
is tested across separated maturities rather than adjacent weeklies. The
surfaces retain between 11 and 20 fitted slices each.

\subsection{Local volatility}

The simulation diffusion is Dupire local volatility, taken in Gatheral's
log-moneyness form,
%
\begin{equation}
  \lv^2(\logm, T) = \frac{\partial\totvar/\partial T}
    {1 - \dfrac{\logm}{\totvar}\dfrac{\partial\totvar}{\partial\logm}
      + \dfrac14\left(-\dfrac14 - \dfrac{1}{\totvar}
        + \dfrac{\logm^2}{\totvar^2}\right)
        \left(\dfrac{\partial\totvar}{\partial\logm}\right)^{\!2}
      + \dfrac12\dfrac{\partial^2\totvar}{\partial\logm^2}}.
  \label{eq:dupire}
\end{equation}
%
Derivatives in $\logm$ are exact. In $T$ the surface is interpolated linearly
in total variance between adjacent expiries, which makes
$\partial\totvar/\partial T$ piecewise constant but non-negative, because the
calendar condition already forced $\totvar$ to be non-decreasing in $T$. The
butterfly condition keeps the denominator of \eqref{eq:dupire} positive. Local
variance is therefore positive by construction rather than by inspection, which
is the reason the two checks in the previous subsection are not optional. The
implementation still counts how often it has to clamp, so a failure would be
visible rather than silent.

Past the longest listed expiry the surface must be extrapolated, and for the
five-year notes that is a long way. Quoted expiries reach 5.20 years on the
S\&P 500 but stop at 3.20 on the Nasdaq-100 and 2.19 on the Russell 2000, while
the two longest notes run to 5.00 years. The base case holds implied volatility
flat in maturity with its smile shape intact,
$\totvar(\logm,T) = \totvar_{\mathrm{last}}(\logm)\,T/T_{\mathrm{last}}$. Two
alternatives are implemented so the choice can be measured rather than
asserted: carrying the last observed variance slope forward, and shifting every
volatility by a fixed amount. \S\ref{sec:inputs} reports what the choice is
worth, which ranges from \$0.22 on the shortest note to \$12.36 on the longest.

\subsection{Correlation}

Correlation is estimated from 252 trading days of daily log returns, taken up to
and including each note's own pricing date rather than the snapshot date. Five
of the six notes priced two business days before the snapshot, so using
everything available would feed future returns into a historical pricing
decision; the effect on the matrix is small and the principle is not. Where the
sample matrix is not positive semidefinite, the nearest valid correlation matrix
in the Frobenius norm is substituted by the alternating-projection method of
\citet{Higham2002}.

This is the largest approximation in the method and it runs in a known
direction. Vanilla options carry no information about co-movement, and there is
no free source of implied correlation, so a realised estimate stands in for a
risk-neutral input. Dealers mark correlation above realised for precisely this
product, because selling a worst-of leaves the desk short correlation. The
investor is long correlation on every note in the sample, so if the issuer used
a higher figure than the one estimated here, its value would exceed this
paper's for that reason alone. That makes correlation the first candidate
\S\ref{sec:inputs} has to rule in or out, and the one it treats at most length.

\subsection{Simulation}

Paths are generated by daily log-Euler steps under the risk-neutral measure,
each asset drifting along its own forward curve so that carry is whatever the
option market implies rather than a separately assumed dividend yield, and
diffusing at \eqref{eq:dupire}. Shocks are correlated through the Cholesky
factor of the correlation matrix. Antithetic variates are used, which help most
where the payoff is discontinuous, and the reported figures use 200{,}000 paths.

Three implementation choices are worth recording because they decide what is
computationally possible rather than merely convenient. Paths are not stored: a
five-year note observed monthly needs 1{,}260 daily steps on three assets, and
retaining 200{,}000 of them would run to several terabytes, so the engine
carries only the current level and records at each observation date the
statistics the payoffs need, which are the worst and best performances across
underlyings. The American knock-in is handled by a running indicator updated
every day, since the condition it tests, that an underlying's value falls below
its threshold ``on any Trading Day during the Knock-In Period,'' cannot be
evaluated from the observation dates alone. That note's loss condition is a
conjunction rather than a single test: a knock-in must occur and the least
performing underlying must also finish below its starting value, not merely
below the threshold that triggered the knock-in. Local volatility is read from a precomputed grid in
$(\logm, T)$ with bilinear interpolation rather than evaluated analytically at
every step, which reproduces the surface to well under a basis point of
volatility and is what makes bump-and-revalue Greeks affordable. Common random
numbers come from re-seeding rather than from stored draws, so a bumped
revaluation differs from the base case only in the bumped input and each Greek
is a clean finite difference.

A European control variate is applied, with the variance reduction factor
running from 1.8 to 3.8 across the six notes. Standard errors on the reported
values fall between \$0.24 and \$0.37 per \$1{,}000.

\subsection{The issuer's call}

The Bermudan call is solved by least-squares Monte Carlo
\citep{LongstaffSchwartz2001}. Treating an issuer's early-redemption right as
an optimal-stopping problem rather than as a barrier test has published
precedent in this product family: \citet{Detemple2021} solve for the optimal
policy explicitly on callable barrier reverse convertibles, and
\citet{Huo2025} applies least-squares Monte Carlo to worst-of issuer-callable
notes of the kind priced here, under a multi-variate local-volatility model and
with the stopping problem posed from the issuer's side. One sign convention has
to be stated plainly, because reversing it would change every result. A high continuation
value is the issuer's liability, not its asset. The issuer therefore redeems
when keeping the note alive costs it more than repaying principal now, so the
investor's value at a call date is the redemption amount whenever the estimated
continuation value exceeds it. Solved the other way the issuer's option becomes
a gift to the investor and every price rises.

The regression basis is deliberately small: a constant, the worst performance
$\worst$, its square, the spread between best and worst performance, and
accrued-but-unpaid coupon. The worst performer is the state variable that drives
the payoff, and its square lets the exercise boundary curve; the best-worst
spread carries the dispersion a worst-of note is exposed to; accrued coupon
matters only where the note has a memory feature. A richer basis fits
simulation noise and estimates a worse boundary rather than a better one.

One detail of the comparison is easy to get wrong in the direction the term
sheets invite. A coupon due on a call date is paid whether or not the note is
redeemed, so it cancels from the decision and the continuation value is
compared against par alone. Term sheets quote the redemption as par plus
coupon, and carrying that wording into the comparison makes calling look about
one coupon more expensive than it is, which systematically under-calls and
overvalues the note.

\subsection{The funding-spread inversion}

The gap between a model value and the disclosed value can be expressed as the
discount-rate adjustment that would remove it. For a constant spread $\spread$
over the Treasury zero curve,
%
\begin{equation}
  V(\spread) = \E^{\mathbb{Q}}\!\left[\;\sum_i
    e^{-(r_{t_i} + \spread)\,t_i}\, \mathrm{CF}_i \right],
  \qquad
  \text{solve } V(\spread^\ast) = \mathrm{EV}_{\text{disclosed}} .
  \label{eq:inversion}
\end{equation}
%
The spread enters the exercise policy as well as the discounting, so
\eqref{eq:inversion} is solved by repricing on a fixed set of paths rather than
by rescaling a single valuation. $V$ is monotone decreasing in $\spread$, so the
root is bracketed and found by Brent's method.

What makes $\spread^\ast$ informative is that its sign is known in advance.
Discounting a liability at a wider spread lowers its value, and an issuer funds
above Treasuries, so a disclosed value derived with the issuer's own funding
rate should reconcile at a positive $\spread^\ast$. A negative solution cannot
be a funding spread. Because the issuer's call is the largest discretionary
input in the model, \eqref{eq:inversion} is solved under two assumptions about
it, optimal exercise and no early call, and \S\ref{sec:results} reports both.
